\documentclass[10pt,a4paper]{amsart}
\usepackage{amsmath,amssymb,amsthm,graphicx}
% Retained minimal preamble for the Stacks Project chapter varieties.tex.
% Verbatim excerpt of the chapter preamble: Xy-pic setup (native preamble lines 14-18), the native theorem
% declarations with the mathematical macros (native preamble lines 154-217), and the project own no-op
% definitions of the reference/slogan/history display environments (native wrapper.tex lines 1-3) that the
% counted statement uses for its slogan. Copied mechanically; nothing added, removed or redefined.
% For commutative diagrams we use Xy-pic
\usepackage[all]{xy}

% We use 2cell for 2-commutative diagrams.
\xyoption{2cell}
% Theorem environments.
%
\theoremstyle{plain}
\newtheorem{theorem}[subsection]{Theorem}
\newtheorem{proposition}[subsection]{Proposition}
\newtheorem{lemma}[subsection]{Lemma}

\theoremstyle{definition}
\newtheorem{definition}[subsection]{Definition}
\newtheorem{example}[subsection]{Example}
\newtheorem{exercise}[subsection]{Exercise}
\newtheorem{situation}[subsection]{Situation}

\theoremstyle{remark}
\newtheorem{remark}[subsection]{Remark}
\newtheorem{remarks}[subsection]{Remarks}

\numberwithin{equation}{subsection}

% Macros
%
\def\lim{\mathop{\mathrm{lim}}\nolimits}
\def\colim{\mathop{\mathrm{colim}}\nolimits}
\def\Spec{\mathop{\mathrm{Spec}}}
\def\Hom{\mathop{\mathrm{Hom}}\nolimits}
\def\Ext{\mathop{\mathrm{Ext}}\nolimits}
\def\SheafHom{\mathop{\mathcal{H}\!\mathit{om}}\nolimits}
\def\SheafExt{\mathop{\mathcal{E}\!\mathit{xt}}\nolimits}
\def\Sch{\mathit{Sch}}
\def\Mor{\mathop{\mathrm{Mor}}\nolimits}
\def\Ob{\mathop{\mathrm{Ob}}\nolimits}
\def\Sh{\mathop{\mathit{Sh}}\nolimits}
\def\NL{\mathop{N\!L}\nolimits}
\def\CH{\mathop{\mathrm{CH}}\nolimits}
\def\proetale{{pro\text{-}\acute{e}tale}}
\def\etale{{\acute{e}tale}}
\def\QCoh{\mathit{QCoh}}
\def\Ker{\mathop{\mathrm{Ker}}}
\def\Im{\mathop{\mathrm{Im}}}
\def\Coker{\mathop{\mathrm{Coker}}}
\def\Coim{\mathop{\mathrm{Coim}}}

% Boxtimes
%
\DeclareMathSymbol{\boxtimes}{\mathbin}{AMSa}{"02}

%
% Macros for moduli stacks/spaces
%
\def\QCohstack{\mathcal{QC}\!\mathit{oh}}
\def\Cohstack{\mathcal{C}\!\mathit{oh}}
\def\Spacesstack{\mathcal{S}\!\mathit{paces}}
\def\Quotfunctor{\mathrm{Quot}}
\def\Hilbfunctor{\mathrm{Hilb}}
\def\Curvesstack{\mathcal{C}\!\mathit{urves}}
\def\Polarizedstack{\mathcal{P}\!\mathit{olarized}}
\def\Complexesstack{\mathcal{C}\!\mathit{omplexes}}
% \Pic is the operator that assigns to X its picard group, usage \Pic(X)
% \Picardstack_{X/B} denotes the Picard stack of X over B
% \Picardfunctor_{X/B} denotes the Picard functor of X over B
\def\Pic{\mathop{\mathrm{Pic}}\nolimits}
\def\Picardstack{\mathcal{P}\!\mathit{ic}}
\def\Picardfunctor{\mathrm{Pic}}
\def\Deformationcategory{\mathcal{D}\!\mathit{ef}}
\newenvironment{reference}{}{}
\newenvironment{slogan}{}{}
\newenvironment{history}{}{}
\usepackage{cleveref}
\crefname{originaldefinition}{original definition}{original definitions}
\Crefname{originaldefinition}{Original definition}{Original definitions}
\makeatletter
\newlabel{morphisms-definition-normalization}{{035N}{source}{}{sourcepointer-1}{}}
\newlabel{morphisms-definition-normalization@cref}{{[originaldefinition][0][]035N}{[1][0][]source}{}{}{}}
\makeatother
\crefname{originallemma}{original lemma}{original lemmas}
\Crefname{originallemma}{Original lemma}{Original lemmas}
\makeatletter
\newlabel{morphisms-lemma-finite-type-noetherian}{{01T6}{source}{}{sourcepointer-2}{}}
\newlabel{morphisms-lemma-finite-type-noetherian@cref}{{[originallemma][0][]01T6}{[1][0][]source}{}{}{}}
\makeatother
\crefname{originallemma}{original lemma}{original lemmas}
\Crefname{originallemma}{Original lemma}{Original lemmas}
\makeatletter
\newlabel{morphisms-lemma-nagata-normalization}{{035S}{source}{}{sourcepointer-3}{}}
\newlabel{morphisms-lemma-nagata-normalization@cref}{{[originallemma][0][]035S}{[1][0][]source}{}{}{}}
\makeatother
\crefname{originallemma}{original lemma}{original lemmas}
\Crefname{originallemma}{Original lemma}{Original lemmas}
\makeatletter
\newlabel{morphisms-lemma-normalization-birational}{{0BXC}{source}{}{sourcepointer-4}{}}
\newlabel{morphisms-lemma-normalization-birational@cref}{{[originallemma][0][]0BXC}{[1][0][]source}{}{}{}}
\makeatother
\crefname{originallemma}{original lemma}{original lemmas}
\Crefname{originallemma}{Original lemma}{Original lemmas}
\makeatletter
\newlabel{morphisms-lemma-normalization-normal}{{035Q}{source}{}{sourcepointer-5}{}}
\newlabel{morphisms-lemma-normalization-normal@cref}{{[originallemma][0][]035Q}{[1][0][]source}{}{}{}}
\makeatother
\crefname{originallemma}{original lemma}{original lemmas}
\Crefname{originallemma}{Original lemma}{Original lemmas}
\makeatletter
\newlabel{morphisms-lemma-normalization-reduced}{{035O}{source}{}{sourcepointer-6}{}}
\newlabel{morphisms-lemma-normalization-reduced@cref}{{[originallemma][0][]035O}{[1][0][]source}{}{}{}}
\makeatother
\crefname{originallemma}{original lemma}{original lemmas}
\Crefname{originallemma}{Original lemma}{Original lemmas}
\makeatletter
\newlabel{morphisms-lemma-ubiquity-nagata}{{035B}{source}{}{sourcepointer-7}{}}
\newlabel{morphisms-lemma-ubiquity-nagata@cref}{{[originallemma][0][]035B}{[1][0][]source}{}{}{}}
\makeatother
\crefname{originallemma}{original lemma}{original lemmas}
\Crefname{originallemma}{Original lemma}{Original lemmas}
\makeatletter
\newlabel{properties-lemma-Noetherian-irreducible-components}{{0BA8}{source}{}{sourcepointer-8}{}}
\newlabel{properties-lemma-Noetherian-irreducible-components@cref}{{[originallemma][0][]0BA8}{[1][0][]source}{}{}{}}
\makeatother
\title{Finiteness, birationality and the reduced case for the normalization of a locally algebraic scheme}
\author{The Stacks Project Authors}
\date{stacks-project:varieties.tex:local-snapshot-20261011; GNU FDL 1.2 or later}
\begin{document}
\maketitle
\noindent\textit{Editorial note.} All delivered mathematical text is the author's own, copied from the pinned Stacks Project chapter file varieties.tex; no statement or proof step was rewritten, reordered or removed.\par\medskip
\noindent\textit{Editorial note.} Parts (2) and (3) are discharged by the author with the single sentences 'Part (2) is Morphisms, Lemma ...' and 'Part (3) is Morphisms, Lemma ...'; these sentences are reproduced verbatim, and they are navigation pointers rather than arguments, exactly as the author wrote them.\par\medskip
\noindent\textit{Editorial note.} The counted spans carry 9 navigation stubs and no bibliography citation.\par\medskip
\noindent\textit{Editorial note.} The preamble is the chapter's retained minimal preamble (sources/stacks-varieties/retained-preamble.tex lines 1-77), a verbatim excerpt of the native preamble.\par\medskip
\noindent\textit{Editorial note.} The label lemma-normalization-locally-algebraic is defined only in varieties.tex, so no cross-chapter restatement had to be weighed for this card.\par\medskip
\noindent\textit{Allowed standard prerequisites (editorial).}\begin{itemize}
\item Locally finite type schemes over a field are locally Noetherian and universally Nagata. Morphisms, Lemmas morphisms-lemma-finite-type-noetherian and morphisms-lemma-ubiquity-nagata, invoked at proof lines 4921-4923 and 4931-4934. Supplies the applicability of the normalization definition and the finiteness of the normalization morphism.
\item Normalization is a disjoint union of normal integral schemes with dominant morphism. Morphisms, Lemma morphisms-lemma-normalization-normal, invoked at proof lines 4928-4930. Supplies the structural part of (1).
\item Normalization of a reduced scheme, and birationality for reduced algebraic schemes. Morphisms, Lemmas morphisms-lemma-normalization-reduced and morphisms-lemma-normalization-birational, invoked at proof lines 4939 and 4942. Supplies parts (2) and (3) as direct references.
\item A quasi-compact open of a locally Noetherian scheme has finitely many irreducible components. Properties, Lemma properties-lemma-Noetherian-irreducible-components, invoked at proof lines 4924-4926. Supplies that the normalization definition applies.
\end{itemize}\medskip
\ifdefined\hypertarget\hypertarget{sourcepointer-1}{}\fi
\noindent\textit{Source reference.} Original definition 035N: sources/stacks-chow/tags.txt, source lines 4134--4134. Pointer outside the counted span, invoked once for context; the load-bearing argument of the counted statement is the author's proof reproduced in full in the proof span.\par\medskip
\ifdefined\hypertarget\hypertarget{sourcepointer-2}{}\fi
\noindent\textit{Source reference.} Original lemma 01T6: sources/stacks-chow/tags.txt, source lines 2363--2363. Pointer outside the counted span, invoked once for context; the load-bearing argument of the counted statement is the author's proof reproduced in full in the proof span.\par\medskip
\ifdefined\hypertarget\hypertarget{sourcepointer-3}{}\fi
\noindent\textit{Source reference.} Original lemma 035S: sources/stacks-chow/tags.txt, source lines 4139--4139. Pointer outside the counted span, invoked once for context; the load-bearing argument of the counted statement is the author's proof reproduced in full in the proof span.\par\medskip
\ifdefined\hypertarget\hypertarget{sourcepointer-4}{}\fi
\noindent\textit{Source reference.} Original lemma 0BXC: sources/stacks-chow/tags.txt, source lines 14975--14975. Pointer outside the counted span, invoked once for context; the load-bearing argument of the counted statement is the author's proof reproduced in full in the proof span.\par\medskip
\ifdefined\hypertarget\hypertarget{sourcepointer-5}{}\fi
\noindent\textit{Source reference.} Original lemma 035Q: sources/stacks-chow/tags.txt, source lines 4137--4137. Pointer outside the counted span, invoked once for context; the load-bearing argument of the counted statement is the author's proof reproduced in full in the proof span.\par\medskip
\ifdefined\hypertarget\hypertarget{sourcepointer-6}{}\fi
\noindent\textit{Source reference.} Original lemma 035O: sources/stacks-chow/tags.txt, source lines 4135--4135. Pointer outside the counted span, invoked once for context; the load-bearing argument of the counted statement is the author's proof reproduced in full in the proof span.\par\medskip
\ifdefined\hypertarget\hypertarget{sourcepointer-7}{}\fi
\noindent\textit{Source reference.} Original lemma 035B: sources/stacks-chow/tags.txt, source lines 4122--4122. Pointer outside the counted span, invoked once for context; the load-bearing argument of the counted statement is the author's proof reproduced in full in the proof span.\par\medskip
\ifdefined\hypertarget\hypertarget{sourcepointer-8}{}\fi
\noindent\textit{Source reference.} Original lemma 0BA8: sources/stacks-chow/tags.txt, source lines 14201--14201. Pointer outside the counted span, invoked once for context; the load-bearing argument of the counted statement is the author's proof reproduced in full in the proof span.\par\medskip
\begin{lemma}
\label{lemma-normalization-locally-algebraic}
Let $k$ be a field. Let $X$ be a locally algebraic scheme over $k$.
Let $\nu : X^\nu \to X$ be the normalization morphism, see
Morphisms, Definition \ref{morphisms-definition-normalization}.
Then
\begin{enumerate}
\item $\nu$ is finite, dominant, and $X^\nu$ is a disjoint
union of normal irreducible locally algebraic schemes over $k$,
\item $\nu$ factors as $X^\nu \to X_{red} \to X$ and the first
morphism is the normalization morphism of $X_{red}$,
\item if $X$ is a reduced algebraic scheme, then $\nu$ is
birational,
\item if $X$ is a variety, then $X^\nu$ is a variety and
$\nu$ is a finite birational morphism of varieties.
\end{enumerate}
\end{lemma}

\begin{proof}
Since $X$ is locally of finite type over a field, we see that
$X$ is locally Noetherian
(Morphisms, Lemma \ref{morphisms-lemma-finite-type-noetherian})
hence every quasi-compact open has finitely many irreducible
components (Properties, Lemma
\ref{properties-lemma-Noetherian-irreducible-components}).
Thus Morphisms, Definition \ref{morphisms-definition-normalization} applies.
The normalization $X^\nu$ is always a disjoint union of normal integral
schemes and the normalization morphism $\nu$ is always dominant, see
Morphisms, Lemma \ref{morphisms-lemma-normalization-normal}.
Since $X$ is universally Nagata
(Morphisms, Lemma \ref{morphisms-lemma-ubiquity-nagata})
we see that $\nu$ is finite
(Morphisms, Lemma \ref{morphisms-lemma-nagata-normalization}).
Hence $X^\nu$ is locally algebraic too.
At this point we have proved (1).

\medskip\noindent
Part (2) is Morphisms, Lemma \ref{morphisms-lemma-normalization-reduced}.

\medskip\noindent
Part (3) is Morphisms, Lemma \ref{morphisms-lemma-normalization-birational}.

\medskip\noindent
Part (4) follows from (1), (2), (3), and the fact that $X^\nu$ is
separated as a scheme finite over a separated scheme.
\end{proof}

\end{document}
